\documentclass[../readings.tex]{subfiles}

\begin{document}

\subsection{Chebyshev Polynomials}
\label{sub:chebyshev-theory}

Minimax-optimal basis for polynomial approximation. Resolves ill-conditioning and oscillations of the monomial basis.

\subsubsection{Failure of the Monomial Basis}

\begin{enumerate}
    \item \textbf{Numerical:}\enspace Vandermonde matrices $V_{ij} = x_i^{j-1}$ have condition numbers that grow exponentially with degree.
    \item \textbf{Analytical:}\enspace \textbf{Runge's Phenomenon}. Interpolating smooth functions at equally spaced nodes leads to divergence at boundaries as degree $n \to \infty$.
\end{enumerate}

\addExcr{%
\begin{enumerate}
    \item Build $n=15$ Vandermonde matrix on equispaced nodes. Check $\kappa_2(\bV)$.
    \item Plot Runge's function $f(x) = 1/(1+25x^2)$ and its $n=10$ interpolant. Note boundary oscillations.
\end{enumerate}
}

\subsubsection{Chebyshev Polynomials ($T_k$)}

Defined on $[-1, 1]$ by $T_k(x) = \cos(k \arccos x)$.

\addDef{Properties of $T_k$}{%
\begin{itemize}
    \item \textbf{Recurrence:}\enspace $T_0=1, T_1=x, T_{k+1}=2x T_k - T_{k-1}$.
    \item \textbf{Boundedness:}\enspace $|T_k(x)| \leq 1$ for all $x \in [-1, 1]$.
    \item \textbf{Orthogonality:}\enspace $\int_{-1}^1 T_j T_k \frac{dx}{\sqrt{1-x^2}} = 0$ for $j \neq k$.
\end{itemize}
}
\label{def:cheb-props}

\addExcr{%
\begin{enumerate}
    \item Compute $T_3, T_4$ via recurrence.
    \item Plot $T_0, \dots, T_5$. Note they all oscillate between $\pm 1$.
\end{enumerate}
}

\subsubsection{Chebyshev Nodes}

Zeros of $T_n(x)$ provide the optimal node distribution for interpolation.

\addDef{Chebyshev Nodes}{%
$x_k = \cos(\frac{(2k-1)\pi}{2n})$ for $k=1, \dots, n$.
\begin{itemize}
    \item \textbf{Node Clustering:}\enspace Nodes cluster near boundaries $\pm 1$. This suppresses Runge's oscillations.
\end{itemize}
}
\label{def:cheb-nodes}

\addNote{%
\textbf{The Minimax Property:}\enspace
Among all monic polynomials of degree $n$, $\tilde{T}_n = T_n/2^{n-1}$ has the smallest possible maximum value on $[-1, 1]$. Using these nodes minimizes the interpolation error bound.
}

\addExcr{%
\begin{enumerate}
    \item Plot 10 Chebyshev nodes vs.\ 10 equispaced nodes. Note clustering.
    \item Repeat the Runge function experiment with Chebyshev nodes. Observe convergence.
    \item Compare $\kappa_2$ of Vandermonde vs.\ Chebyshev basis ($V_{ij} = T_{j-1}(x_i)$).
\end{enumerate}
}

\subsubsection{Chebyshev Series}

Truncated expansions $f(x) \approx \sum c_k T_k(x)$.

\addThm{Spectral Accuracy}{%
If $f$ is analytic on $[-1, 1]$, the Chebyshev coefficients $c_k$ decay \textbf{exponentially}. Truncation error $\|f - p_n\|_\infty = O(\rho^{-n})$.
\textbf{Insight:}\enspace Just 20--30 terms often yield machine precision.
}
\label{thm:spectral-accuracy}

\addExcr{%
\begin{enumerate}
    \item Compute degree-$n$ fit to $e^x$ via \texttt{np.polynomial.chebyshev.chebfit}. Plot error vs.\ $n$.
    \item \textbf{Spectral Differentiation:}\enspace The derivative of a Chebyshev series satisfies a simple recurrence. Compare \texttt{chebder} accuracy to finite differences for $f(x) = \sin(\pi x)$.
\end{enumerate}
}

\end{document}
